\chapter{Rodziny sprzężone: rachunek posterioru i prognoza}
\label{ch:conjugate-families}
\paragraph{Powiązania i odsyłacze.} Wzór Bayesa i warunki normalizacji wprowadziliśmy w rozdziale~\ref{ch:bayes-foundations}; predykcję rozwijamy w rozdziale~\ref{ch:bayes-estimation-prediction}.


\section{Dlaczego sprzężoność jest użyteczna?}
Sprzężoność pozwala przekształcić wzór Bayesa w jawny rachunek:
po obserwacji danych rozkład a posteriori pozostaje w tej samej
rodzinie co prior. W praktyce otrzymujemy dzięki temu proste
prognozy, przedziały wiarygodności i wygodne aktualizowanie wiedzy
po napływie nowych obserwacji. Nie jest to jednak argument za
wyborem rodziny jako \emph{prawdziwego} modelu danych:
wygoda obliczeniowa nie zastępuje diagnostyki ani analizy wrażliwości.

\begin{definicja}[Rodzina sprzężona]
Dla jądra wiarygodności $p(x\mid\theta)$ na
$\Theta_0$ rodzinę \emph{właściwych} priorów
$\{\pi_\eta:\eta\in H\}$ nazywamy sprzężoną, jeśli dla
każdego dopuszczalnego $\eta$ i każdego $x$ o dodatniej,
skończonej gęstości predykcyjnej rozkład a posteriori jest
równy $\pi_{T(\eta,x)}$ dla pewnego odwzorowania aktualizacji
$T$. Sprzężoność określa parę (model wiarygodności, prior),
a nie własność samej rodziny priorów.
\end{definicja}

\section{Beta--Bernoulli i beta--dwumianowy: szacowanie odsetka}
\begin{definicja}[Rozkład beta]
Dla $a,b>0$ rozkład $\operatorname{Beta}(a,b)$ na $(0,1)$
ma gęstość
\[
 f_{a,b}(p)=\frac{p^{a-1}(1-p)^{b-1}}{B(a,b)},
 \qquad B(a,b)=\int_0^1u^{a-1}(1-u)^{b-1}\,\dd\lambda_1(u)
 =\frac{\Gamma(a)\Gamma(b)}{\Gamma(a+b)}.
\]
Funkcja $\Gamma(t)=\int_0^\infty u^{t-1}e^{-u}\,\dd\lambda_1(u)$
dla $t>0$ jest funkcją gamma. Tożsamość beta--gamma wynika
ze zmiany zmiennych w podwójnej całce gamma; wykorzystamy
ją do obliczania momentów.
\end{definicja}

\begin{lemat}[Tożsamość beta--gamma]
Dla wszystkich $a,b>0$ zachodzi
$B(a,b)=\Gamma(a)\Gamma(b)/\Gamma(a+b)$.
\end{lemat}
\begin{proof}
Korzystając z twierdzenia Tonellego, zapisujemy iloczyn
$\Gamma(a)\Gamma(b)$ jako całkę funkcji
$x^{a-1}y^{b-1}e^{-x-y}$ po $(0,\infty)^2$
względem $\lambda_2$. Podstawiamy $x=ru$, $y=r(1-u)$
dla $r>0$, $u\in(0,1)$. Jakobian ma wartość $r$,
więc nowa całka jest iloczynem
$\int_0^\infty r^{a+b-1}e^{-r}\,\dd\lambda_1(r)$
i $\int_0^1u^{a-1}(1-u)^{b-1}\,\dd\lambda_1(u)$.
Pierwsza całka równa się $\Gamma(a+b)$, druga $B(a,b)$.
\end{proof}

\begin{twierdzenie}[Sprzężoność i rozkład predykcyjny beta--dwumianowy]
\label{thm:beta-binomial}
Niech $p\sim\operatorname{Beta}(a,b)$, $a,b>0$, a warunkowo
względem $p$ zmienne $X_1,\ldots,X_n$ będą niezależne
i $X_i\mid p\sim\operatorname{Bernoulli}(p)$, gdzie $n\ge1$.
Niech $S=\sum_{i=1}^nX_i$ i $s\in\{0,\ldots,n\}$.
Wtedy
\[
 p\mid S=s\sim\operatorname{Beta}(a+s,b+n-s),\qquad
 \E[p\mid S=s]=\frac{a+s}{a+b+n},
\]
\[
 \Var(p\mid S=s)=\frac{(a+s)(b+n-s)}{(a+b+n)^2(a+b+n+1)}.
\]
Jeżeli po próbie otrzymamy $m\ge1$ nowych, warunkowo niezależnych
obserwacji Bernoulliego o tym samym parametrze i
$T$ oznacza liczbę sukcesów, to dla $j=0,\ldots,m$
\[
 \Prob(T=j\mid S=s)
 =\binom mj\frac{B(a+s+j,b+n-s+m-j)}{B(a+s,b+n-s)}.
\]
W szczególności przyszły pojedynczy sukces ma prawdopodobieństwo
$(a+s)/(a+b+n)$.
\end{twierdzenie}
\begin{proof}
Prawdopodobieństwo $S=s$ przy znanym $p$ wynosi
$\binom ns p^s(1-p)^{n-s}$. Iloczyn gęstości prioru i
wiarygodności jest proporcjonalny do
$p^{a+s-1}(1-p)^{b+n-s-1}$, więc normalizacja przez
$B(a+s,b+n-s)$ daje rozkład wskazany w tezie.
Dla $r>0$, $t>0$, $k\in\{1,2\}$, z definicji beta
$\E[U^k]=B(r+k,t)/B(r,t)$, gdzie
$U\sim\operatorname{Beta}(r,t)$. Tożsamość
$\Gamma(z+1)=z\Gamma(z)$ daje
$\E U=r/(r+t)$ oraz $\Var(U)=rt/[(r+t)^2(r+t+1)]$.
Po wstawieniu $r=a+s$, $t=b+n-s$ otrzymujemy wzory
momentowe. Dla nowych obserwacji
$\Prob(T=j\mid p)=\binom mjp^j(1-p)^{m-j}$.
Całkując to prawdopodobieństwo względem wyprowadzonego
posterioru, otrzymujemy ostatni wzór.
\end{proof}

\begin{stwierdzenie}[Średnia posterioru jako kompromis MLE i prioru]
\label{prop:beta-shrinkage}
W założeniach twierdzenia~\ref{thm:beta-binomial} oznaczmy
$\widehat p=S/n$, $m_0=a/(a+b)$ i $w=n/(n+a+b)$.
Wówczas $\E[p\mid S]=w\widehat p+(1-w)m_0$.
\end{stwierdzenie}
\begin{proof}
Po podstawieniu dostajemy
$\frac{n}{n+a+b}\frac Sn+
\frac{a+b}{n+a+b}\frac{a}{a+b}=(S+a)/(n+a+b)$.
\end{proof}

\begin{przyklad}[Kontrola jakości: 7 wad na 100 sztuk]
Dla $S=7$, $n=100$ estymator MLE wynosi $0{,}07$.
Z priorem $\operatorname{Beta}(1,1)$ posterior ma rozkład
$\operatorname{Beta}(8,94)$, średnią $8/102\approx0{,}0784$,
a prawdopodobieństwo wady następnej sztuki wynosi także $8/102$.
Dla prioru $\operatorname{Beta}(20,80)$ posterior ma rozkład
$\operatorname{Beta}(27,173)$ i średnią $0{,}135$.
Obydwa modele używają identycznych danych, lecz różnią się
założeniami; prior o łącznej masie parametrów $a+b=100$
ma wyraźny wpływ na wynik. „Siła prioru” $a+b$ jest tu
wygodnym skrótem algebraicznym, nie dokładną liczbą
rzeczywistych dodatkowych niezależnych obserwacji.
Przedział wiarygodności wyznaczymy z kwantyli rozkładu beta,
a jego interpretację odróżnimy od przedziału Wilsona z
rozdziału~\ref{ch:confidence-intervals}.
\end{przyklad}

\begin{uwaga}[MAP i przypadki na brzegu]
Dla posterioru $\operatorname{Beta}(r,t)$, $r,t>1$,
moda leży w $(r-1)/(r+t-2)$. Jeśli $r\le1$ lub $t\le1$,
maksimum gęstości może wypadać na brzegu lub nie być osiągane
na otwartym $(0,1)$. Gdy $r=t=1$, gęstość jest stała:
każdy punkt $(0,1)$ jest jej dominantą. Przy $r<1$ lub
$t<1$ gęstość może być nieograniczona w pobliżu brzegu,
a wówczas \emph{maksymalizator w otwartym przedziale nie istnieje}.
Nie wolno bezwarunkowo używać wzoru na modę dla dowolnych
$r,t>0$. Sam MAP nie jest niezmienniczy względem nieliniowej
reparametryzacji; podobny problem dotyczy definicji HPD.
\end{uwaga}

\section{Gamma--Poisson: częstość awarii i ekspozycja}
\begin{definicja}[Rozkład gamma w parametryzacji shape--rate]
Dla $a,b>0$ przez $\operatorname{Gamma}(a,b)$ rozumiemy
rozkład na $(0,\infty)$ o gęstości
$g_{a,b}(\lambda)=b^a\lambda^{a-1}e^{-b\lambda}/\Gamma(a)$.
Pierwszy parametr to \emph{kształt} (shape), drugi to
\emph{intensywność zaniku} (rate), a nie skala.
Średnia wynosi $a/b$, a wariancja $a/b^2$.
Jeżeli $\lambda$ jest liczbą awarii na godzinę, wtedy
$b$ ma jednostkę godzin, tak aby iloczyn $b\lambda$
był bezwymiarowy.
\end{definicja}

\begin{twierdzenie}[Aktualizacja i predykcja gamma--Poisson]
\label{thm:gamma-poisson}
Niech $\lambda\sim\operatorname{Gamma}(a,b)$,
$a,b>0$. Warunkowo przy $\lambda$ obserwujemy niezależne
zliczenia $Y_i\sim\operatorname{Pois}(e_i\lambda)$,
gdzie ekspozycje $e_i>0$ są znanymi liczbami,
$i=1,\ldots,n$, $n\ge1$. Niech $S=\sum_{i=1}^nY_i$
i $E=\sum_{i=1}^ne_i$. Wtedy
$\lambda\mid Y_1,\ldots,Y_n\sim\operatorname{Gamma}(a+S,b+E)$.
Dla nowej ekspozycji $e_*>0$ i nowego zliczenia
$Y_*\mid\lambda\sim\operatorname{Pois}(e_*\lambda)$,
warunkowo niezależnego od danych, zachodzi dla $k\in\N\cup\{0\}$
\[
 \Prob(Y_*=k\mid Y_1,\ldots,Y_n)
 =\frac{\Gamma(a+S+k)}{\Gamma(a+S)k!}
 \left(\frac{b+E}{b+E+e_*}\right)^{a+S}
 \left(\frac{e_*}{b+E+e_*}\right)^k.
\]
Ponadto $\E[Y_*\mid Y]=e_*(a+S)/(b+E)$ i
$\Var(Y_*\mid Y)=e_*(a+S)/(b+E)
+e_*^2(a+S)/(b+E)^2$.
\end{twierdzenie}
\begin{proof}
Warunkowa wiarygodność zliczeń, z dokładnością do
czynnika niezależnego od $\lambda$, wynosi
$\lambda^Se^{-E\lambda}$. Mnożenie przez prior
$\lambda^{a-1}e^{-b\lambda}$ daje jądro gamma
$\lambda^{a+S-1}e^{-(b+E)\lambda}$, które jest
normalizowalne, gdy $a+S,b+E>0$.
Piszmy $r=a+S$, $v=b+E$. Warunkowe prawdopodobieństwo
$Y_*=k$ ma postać
\[
 \frac{e_*^kv^r}{k!\Gamma(r)}
 \int_0^\infty\lambda^{k+r-1}
 e^{-(v+e_*)\lambda}\,\dd\lambda_1(\lambda).
\]
Z definicji funkcji gamma całka wynosi
$\Gamma(r+k)/(v+e_*)^{r+k}$, co daje wzór.
Wreszcie korzystamy z tożsamości
$\E[Y_*\mid\lambda,Y]=e_*\lambda$ i
$\Var(Y_*\mid\lambda,Y)=e_*\lambda$;
wzór na wariancję całkowitą z rozdziału~\ref{ch:bayes-foundations}
daje podane dwa składniki.
\end{proof}

\begin{przyklad}[Awaryjność z różną długością obserwacji]
Obserwujemy 4 awarie podczas 200 godzin działania urządzenia.
Model zliczeń ma ekspozycję $E=200$ godzin i nieznaną
intensywność $\lambda$ na godzinę. Niech prior to
$\operatorname{Gamma}(2,100)$ (shape--rate).
Posterior to $\operatorname{Gamma}(6,300)$ i
$\E[\lambda\mid\mathrm{dane}]=6/300=0{,}02$ na godzinę.
Prognozowana liczba awarii przez następne 50 godzin ma średnią
$50\cdot0{,}02=1$. Jej wariancja wynosi
$1+50^2\cdot 6/300^2=1+1/6$; to więcej niż
wariancja rozkładu Poissona z ustalonym parametrem $1$.
Powodem jest niepewność o $\lambda$.
W praktyce niezależność przyrostów i stała intensywność
w czasie wymagają sprawdzenia; przy sezonowości lub awariach
w klastrach rozważamy model bardziej złożony.
\end{przyklad}

\section{Gamma--wykładniczy: czas oczekiwania}
\begin{twierdzenie}[Posterior intensywności rozkładu wykładniczego]
Niech $\lambda\sim\operatorname{Gamma}(a,b)$, $a,b>0$,
a przy znanym $\lambda$ obserwacje $T_1,\ldots,T_n$
będą niezależne o gęstości
$f(t\mid\lambda)=\lambda e^{-\lambda t}\ind{t\ge0}$
względem miary Lebesgue'a, $n\ge1$.
Dla $s=\sum_{i=1}^n t_i$ posterior ma rozkład
$\operatorname{Gamma}(a+n,b+s)$.
Jeżeli przyszły czas $T_*$ jest warunkowo niezależny od danych,
ma warunkową gęstość wykładniczą o tej samej intensywności,
to dla $u\ge0$
\[
 \Prob(T_*>u\mid T_1=t_1,\ldots,T_n=t_n)
 =\left(\frac{b+s}{b+s+u}\right)^{a+n}.
\]
\end{twierdzenie}
\begin{proof}
Iloczyn gęstości obserwacji daje $\lambda^n e^{-s\lambda}$,
więc posterior jest gamma przez mnożenie przez jego prior.
Po warunkowaniu na $\lambda$ prawdopodobieństwo przetrwania
to $e^{-u\lambda}$. Całkując je względem posterioru
$\operatorname{Gamma}(a+n,b+s)$ i wykorzystując wzór na
funkcję gamma, otrzymujemy wskazaną potęgę.
\end{proof}

\begin{uwaga}[Cenzorowanie a kompletne czasy]
Jeżeli niektórych czasów nie zaobserwowano do końca, zwykły
iloczyn $\prod_i\lambda e^{-\lambda t_i}$ jest błędny.
Przy prawostronnym, niezależnym cenzorowaniu z liczbą
rzeczywiście zaobserwowanych zdarzeń $d$ i łącznym czasem
ekspozycji $E$, część wiarygodności zależna od $\lambda$
ma postać $\lambda^d e^{-\lambda E}$. Posterior
$\operatorname{Gamma}(a+d,b+E)$ pozostaje sprzężony,
ale tylko przy prawidłowym opisaniu mechanizmu cenzorowania.
\end{uwaga}

\section{Normalny model średniej ze znaną wariancją}
\begin{twierdzenie}[Normalny prior dla normalnej średniej]
\label{thm:normal-normal}
Niech $\mu\sim\mathcal N(m_0,v_0)$, $v_0>0$, a warunkowo
względem $\mu$ zmienne $X_1,\ldots,X_n$ będą niezależne
$\mathcal N(\mu,\sigma^2)$ dla znanego $\sigma^2>0$, $n\ge1$.
Wtedy
\[
 \mu\mid X_1=x_1,\ldots,X_n=x_n\sim\mathcal N(m_n,v_n),
 \quad v_n=\left(\frac1{v_0}+\frac n{\sigma^2}\right)^{-1},
 \quad m_n=v_n\left(\frac{m_0}{v_0}+\frac{n\overline x}{\sigma^2}\right).
\]
Średnia posterioru jest średnią ważoną priorowej średniej
$m_0$ i średniej z próby $\overline x$, z wagami
proporcjonalnymi do ich precyzji (odwrotności wariancji).
\end{twierdzenie}
\begin{proof}
Zależność od $\mu$ w priorze i wiarygodności ma postać
$\exp\{-\tfrac12[(\mu-m_0)^2/v_0+
\sum_{i=1}^n(x_i-\mu)^2/\sigma^2]\}$.
Po rozpisaniu kwadratów współczynnik przy $\mu^2$
wynosi $v_0^{-1}+n\sigma^{-2}=v_n^{-1}$,
a współczynnik przy $-2\mu$ jest
$m_0/v_0+n\overline x/\sigma^2=v_n^{-1}m_n$.
Uzupełnienie kwadratu daje jądro rozkładu
$\mathcal N(m_n,v_n)$; normalizacja jest skończona.
\end{proof}

\begin{przyklad}[Pomiar z kalibracją a nowe obserwacje]
Z wcześniejszej kalibracji wynika $\mu\sim\mathcal N(10,1)$.
Pięć nowych niezależnych pomiarów ma znaną wariancję
błędu $\sigma^2=4$ i średnią $12$. Posteriorowa precyzja
wynosi $1+5/4=2{,}25$, zatem
$v_n=4/9\approx0{,}4444$ i
$m_n=(10+5\cdot12/4)/(1+5/4)=100/9\approx11{,}111$.
Wartość średniej została przesunięta od priorowego $10$ ku
obserwowanemu $12$. Jeżeli wariancja pomiaru jest nieznana,
ten rozkład przestaje być automatycznie właściwym modelem.
\end{przyklad}

\section{Nieznane średnia i wariancja: normalny--odwrotny gamma}
\begin{definicja}[Rozkład odwrotny gamma]
Dla $\alpha,\beta>0$ rozkład
$\operatorname{InvGamma}(\alpha,\beta)$ na $(0,\infty)$
ma gęstość
$h(v)=\beta^\alpha v^{-\alpha-1}e^{-\beta/v}/\Gamma(\alpha)$.
Parametr $\beta$ jest parametrem \emph{skali w wykładniku},
a nie intensywnością gamma na $v$. Jeżeli
$V\sim\operatorname{InvGamma}(\alpha,\beta)$,
to $V^{-1}\sim\operatorname{Gamma}(\alpha,\beta)$
w parametryzacji shape--rate.
\end{definicja}

\begin{twierdzenie}[Sprzężony model normalny z nieznanymi parametrami]
\label{thm:normal-invgamma}
Niech $X_1,\ldots,X_n$ będą warunkowo niezależne
$\mathcal N(\mu,v)$, gdzie $\mu\in\R$, $v>0$, $n\ge1$.
Niech prior ma postać
$v\sim\operatorname{InvGamma}(\alpha_0,\beta_0)$,
$\mu\mid v\sim\mathcal N(m_0,v/\kappa_0)$,
przy $\alpha_0,\beta_0,\kappa_0>0$, $m_0\in\R$.
Piszmy $\overline x=n^{-1}\sum_ix_i$,
$S=\sum_i(x_i-\overline x)^2$. Wtedy posterior ma tę samą
postać z parametrami
\[
 \kappa_n=\kappa_0+n,\quad
 m_n=\frac{\kappa_0m_0+n\overline x}{\kappa_0+n},\quad
 \alpha_n=\alpha_0+\frac n2,
\]
\[
 \beta_n=\beta_0+\frac S2+
 \frac{\kappa_0n}{2(\kappa_0+n)}(\overline x-m_0)^2.
\]
Dokładniej $v\mid x\sim\operatorname{InvGamma}(\alpha_n,\beta_n)$
i $\mu\mid v,x\sim\mathcal N(m_n,v/\kappa_n)$.
\end{twierdzenie}
\begin{proof}
Wspólny prior ma gęstość względem miary Lebesgue'a $\lambda_2$
na $\R\times(0,\infty)$
proporcjonalną do
$v^{-\alpha_0-3/2}\exp\{-\beta_0/v-
\kappa_0(\mu-m_0)^2/(2v)\}$.
Wiarygodność jest proporcjonalna do
$v^{-n/2}\exp\{-\sum_i(x_i-\mu)^2/(2v)\}$.
Korzystamy z dwóch tożsamości:
$\sum_i(x_i-\mu)^2=S+n(\mu-\overline x)^2$ oraz
\[
 \kappa_0(\mu-m_0)^2+n(\mu-\overline x)^2
 =\kappa_n(\mu-m_n)^2+
 \frac{\kappa_0n}{\kappa_n}(\overline x-m_0)^2.
\]
Po podstawieniu wspólne jądro posterioru jest równe
$v^{-\alpha_n-3/2}\exp\{-\beta_n/v-
\kappa_n(\mu-m_n)^2/(2v)\}$.
Jest to jądro gęstości normalnej warunkowej pomnożonej przez
gęstość odwrotną gamma z parametrami $\alpha_n,\beta_n$;
obydwa czynniki są normalizowalne.
\end{proof}

\begin{wniosek}[Predykcja t-Studenta]
\label{cor:student-predictive}
W założeniach twierdzenia~\ref{thm:normal-invgamma}
niech $X_*$ będzie nową obserwacją warunkowo niezależną
od próby przy znanych $\mu,v$ i niech $X_*\mid\mu,v\sim
\mathcal N(\mu,v)$. Wtedy $X_*\mid x$ ma rozkład
$t$-Studenta z $2\alpha_n$ stopniami swobody, środkiem $m_n$
i parametrem skali
$s_n=\sqrt{\beta_n(\kappa_n+1)/(\alpha_n\kappa_n)}$,
w tej konwencji, że $(X_*-m_n)/s_n$ ma standardowy rozkład
$t_{2\alpha_n}$.
\end{wniosek}
\begin{proof}
Przy ustalonym $v$ całkujemy najpierw po $\mu$;
w wyniku splotu normalnego dostajemy
$X_*\mid v,x\sim\mathcal N(m_n,v(1+\kappa_n^{-1}))$.
Zatem przy niezależnych $Z\sim\mathcal N(0,1)$
i $G\sim\operatorname{Gamma}(\alpha_n,\beta_n)$
można przedstawić $X_*-m_n=\sqrt{(1+\kappa_n^{-1})/G}\,Z$.
Ponieważ $2\beta_n G\sim\chi^2_{2\alpha_n}$
(z własności transformacji rozkładu gamma), stosunek
$Z/\sqrt{2\beta_nG/(2\alpha_n)}$ ma standardowy rozkład
$t_{2\alpha_n}$. Po przeliczeniu współczynnika otrzymujemy
$s_n^2=\beta_n(\kappa_n+1)/(\alpha_n\kappa_n)$.
\end{proof}

\begin{uwaga}[Uwaga na istnienie momentów]
Dla posterioru $\operatorname{InvGamma}(\alpha_n,\beta_n)$
wartość oczekiwana wariancji $v$ istnieje dopiero przy
$\alpha_n>1$ i wynosi $\beta_n/(\alpha_n-1)$.
Wariancja rozkładu predykcyjnego t-Studenta jest skończona
dopiero przy $2\alpha_n>2$. Nie należy podawać nieistniejących
momentów w modelach z bardzo słabą informacją.
\end{uwaga}

\section{Dirichlet--wielomianowy: więcej niż dwie kategorie}
\begin{definicja}[Rozkład Dirichleta]
Dla $K\ge2$ i $a_1,\ldots,a_K>0$ niech
$\Delta_{K-1}=\{(p_1,\ldots,p_K)\in(0,1)^K:\sum_{j=1}^Kp_j=1\}$.
Rozkład $\operatorname{Dirichlet}(a_1,\ldots,a_K)$ ma
gęstość względem $\dd\lambda_{K-1}(p_1,\ldots,p_{K-1})$
na obszarze $p_j>0$, $\sum_{j<K}p_j<1$, po podstawieniu
$p_K=1-\sum_{j<K}p_j$, proporcjonalną do
$\prod_{j=1}^Kp_j^{a_j-1}$. Stała normalizacji jest
$\Gamma(a_\bullet)/\prod_{j=1}^K\Gamma(a_j)$,
gdzie $a_\bullet=\sum_{j=1}^Ka_j$.
\end{definicja}

\begin{lemat}[Stała normalizująca rozkładu Dirichleta]
Dla $K\ge2$ oraz $a_1,\ldots,a_K>0$ zachodzi
\[
 \int_{D_{K-1}}\prod_{j=1}^{K}p_j^{a_j-1}
 \,\dd\lambda_{K-1}(p_1,\ldots,p_{K-1})
 =\frac{\prod_{j=1}^K\Gamma(a_j)}{\Gamma(a_\bullet)},
\]
gdzie $D_{K-1}=\{(p_1,\ldots,p_{K-1}):p_j>0,
\ \sum_{j<K}p_j<1\}$, $p_K=1-\sum_{j<K}p_j$ oraz
$a_\bullet=\sum_{j=1}^Ka_j$.
\end{lemat}
\begin{proof}
Iloczyn $\prod_{j=1}^K\Gamma(a_j)$ jest całką funkcji
$\prod_{j=1}^Kt_j^{a_j-1}e^{-t_j}$ po $(0,\infty)^K$
względem $\lambda_K$. Wprowadzamy
$r=\sum_{j=1}^Kt_j$ i $p_j=t_j/r$ dla $j<K$;
wtedy $t_K=r(1-\sum_{j<K}p_j)$.
Jakobian tej zamiany zmiennych ma wartość $r^{K-1}$.
Po zastosowaniu Tonellego całka rozdziela się na
$\int_0^\infty r^{a_\bullet-1}e^{-r}\,\dd\lambda_1(r)
=\Gamma(a_\bullet)$ i całkę z tezy. Dzielenie
przez dodatnią wartość $\Gamma(a_\bullet)$ kończy dowód.
\end{proof}

\begin{twierdzenie}[Sprzężoność dla zliczeń wielokategorialnych]
Niech $p=(p_1,\ldots,p_K)$ ma rozkład
$\operatorname{Dirichlet}(a_1,\ldots,a_K)$, a przy znanym $p$
obserwujemy $n\ge1$ niezależnych wyników z $K$ kategorii,
przy czym prawdopodobieństwo kategorii $j$ wynosi $p_j$.
Jeśli w danych kategoria $j$ wystąpiła $N_j$ razy,
$\sum_{j=1}^KN_j=n$, to posterior ma rozkład
$\operatorname{Dirichlet}(a_1+N_1,\ldots,a_K+N_K)$.
Prawdopodobieństwo, że następny wynik trafi do kategorii $j$,
wynosi $(a_j+N_j)/(a_\bullet+n)$.
\end{twierdzenie}
\begin{proof}
Wiarygodność jako funkcja $p$ jest proporcjonalna do
$\prod_{j=1}^Kp_j^{N_j}$. Po pomnożeniu przez gęstość
Dirichleta otrzymujemy jądro
$\prod_{j=1}^Kp_j^{a_j+N_j-1}$, którego normalizacja
jest gęstością posterioru. Średnią współrzędnej
$p_j$ obliczamy przez iloraz stałych normalizujących
$\Gamma(a_j+N_j+1)/\Gamma(a_j+N_j)$ oraz
$\Gamma(a_\bullet+n)/\Gamma(a_\bullet+n+1)$.
\end{proof}

\section{Obliczenia numeryczne i porównanie z metodami klasycznymi}
Posterior warunkowy opisuje parametry \emph{po zobaczeniu danych},
ale jeszcze przed ich obserwacją także prior nakłada konkretne
ograniczenia na możliwą próbę. W rozdziale~\ref{ch:bayes-estimation-prediction}
wyprowadzimy predykcyjny rozkład liczby sukcesów z modelu
beta--dwumianowego i pokażemy, jak może on służyć do oceny
realizmu założeń a priori. Nie wolno zastępować tej kontroli
stwierdzeniem, że algebraiczna sprzężoność jest wygodna.

W obliczeniach posteriorów korzystamy ze stabilnych funkcji
bibliotecznych kwantyli beta, gamma i t-Studenta; bezpośrednie
mnożenie dużej liczby gęstości prowadzi do niedomiaru
zmiennoprzecinkowego. Poniższe programy odtwarzają
parametry posterioru i predykcję, ale nie zastępują
diagnostyki poprawności modelu.

\begin{lstlisting}[language=Python,caption={Kompletny program: beta--dwumianowy, przedziały i test wrażliwości na prior (Python).}]
from scipy.stats import beta

n, defects = 100, 7
for a, b in [(1, 1), (2, 18), (20, 80)]:
    post_a, post_b = a + defects, b + n - defects
    mean = post_a / (post_a + post_b)
    lo, hi = beta.ppf([0.025, 0.975], post_a, post_b)
    prob_over_10pct = beta.sf(0.10, post_a, post_b)
    print((a, b), "mean:", round(mean, 5),
          "95% credible:", (round(lo, 5), round(hi, 5)),
          "P(p>0.10|data):", round(prob_over_10pct, 5))

# Probability of zero defects in the next ten devices.
a1, b1 = 1 + defects, 1 + n - defects
from scipy.special import betaln
import numpy as np
prob_zero = np.exp(betaln(a1, b1 + 10) - betaln(a1, b1))
print("predictive P(0 defects among next 10):", prob_zero)
\end{lstlisting}

\begin{lstlisting}[language=Julia,caption={Kompletny program: aktualizacja sprzężona i predykcja gamma--Poisson (Julia, biblioteka standardowa).}]
using Printf

function gamma_poisson_update(a::Real, b::Real,
                              counts::AbstractVector{<:Integer},
                              exposures::AbstractVector{<:Real})
    length(counts) == length(exposures) ||
        throw(DimensionMismatch("different lengths"))
    a > 0 && b > 0 || throw(ArgumentError("prior must be proper"))
    all(x -> x >= 0, counts) || throw(ArgumentError("negative count"))
    all(x -> x > 0, exposures) || throw(ArgumentError("exposure <= 0"))
    a_new = a + sum(counts)
    b_new = b + sum(exposures)
    return a_new, b_new
end

a, b = gamma_poisson_update(2.0, 100.0, [1, 0, 3],
                            [50.0, 50.0, 100.0])
e_next = 50.0
mean_next = e_next * a / b
var_next = mean_next + e_next^2 * a / b^2
@printf("posterior Gamma(shape=%.1f, rate=%.1f)\n", a, b)
@printf("predictive mean=%.6f, variance=%.6f\n",
        mean_next, var_next)
\end{lstlisting}

\begin{center}\small
\begin{tabular}{>{\raggedright\arraybackslash}p{.26\textwidth}>{\raggedright\arraybackslash}p{.29\textwidth}>{\raggedright\arraybackslash}p{.32\textwidth}}
\toprule
\textbf{Problem / dane} & \textbf{Klasyczne narzędzie} & \textbf{Sprzężony model bayesowski i uwaga}\\
\midrule
Odsetek wad, mała próba & MLE, Wilson, Clopper--Pearson & Beta--Bernoulli; posterior zależy od prioru.\\
Zliczenia z ekspozycją & MLE Poissona, dokładny przedział & Gamma--Poisson; łatwa predykcja z nadmierną dyspersją.\\
Średnia przy znanym $\sigma$ & Średnia, przedział normalny & Normalny--normalny; precyzyjna aktualizacja.\\
Średnia i wariancja przy normalności & $t$-Studenta, $\chi^2$ & Normalny--odwrotny gamma; predykcja t.\\
Wiele kategorii & Udziały empiryczne, wielomianowy & Dirichlet--wielomianowy; prosty posterior.\\
Nieustalony mechanizm zależności & Model czasowy lub hierarchiczny & Sama sprzężoność iid nie usuwa zależności.\\
\bottomrule
\end{tabular}
\end{center}

\begin{zadanie}[Brak awarii nie oznacza zerowej intensywności]
Podczas 50 godzin działania nie zanotowano awarii. Przy priorze
$\lambda\sim\operatorname{Gamma}(1,50)$ i modelu Poissona
z intensywnością $\lambda$ na godzinę wyznacz posterior oraz
prawdopodobieństwo braku awarii w ciągu kolejnych 50 godzin.
\par\emph{Rozwiązanie.} W twierdzeniu~\ref{thm:gamma-poisson}
$S=0$, $E=50$, więc posterior to $\operatorname{Gamma}(1,100)$.
Prawdopodobieństwo $Y_*=0$ dla $e_*=50$ wynosi
$(100/(100+50))^1=2/3$.
Estymator MLE intensywności wynosi $0$, ale nie wynika z tego
zerowa predykcyjna szansa przyszłej awarii.
\end{zadanie}

\begin{zadanie}[Dane i wrażliwość na założenia]
Dla danych $S=7$, $n=100$ rozważ priory
$\operatorname{Beta}(1,1)$ i $\operatorname{Beta}(20,80)$.
Porównaj średnie posteriorów i zinterpretuj różnicę.
\par\emph{Rozwiązanie.} Otrzymujemy odpowiednio
$8/102\approx0{,}0784$ i $27/200=0{,}135$.
Duża różnica przy stałych danych wynika z innej ilości i treści
informacji zakodowanej w priorach. W praktyce należy opisać
pochodzenie prioru i pokazać analizę wrażliwości; nie wystarczy
wybrać prioru, który daje oczekiwany wynik.
\end{zadanie}


% ================================================================
% ROZDZIAL 19 -- ESTYMACJA I PREDYKCJA BAYESOWSKA
% ================================================================
